\documentclass[12pt]{article} 
\usepackage[utf8]{inputenc} 

\usepackage[a4paper,top=1cm,bottom=0.7cm,left=2cm,right=2cm]{geometry} 
\usepackage{amsmath}
\usepackage{amssymb}
\usepackage{graphicx} 
\usepackage{setspace} 
\usepackage{enumitem} 
\usepackage{appendix} 
\usepackage{listings}
\usepackage{xcolor} 
\usepackage{hyperref}
\usepackage{rs-notes}
% This legacy sheet uses fully project-qualified references only.
% Bridge them to the same named destinations used by the course packages.
\newcommand{\xref}[2][]{%
  \begingroup
  \edef\cheatRefTarget{#2}%
  \expandafter\cheatResolveRef\cheatRefTarget\cheatEndRef{#1}%
  \endgroup}
\def\cheatResolveRef#1::#2::#3:#4\cheatEndRef#5{%
  \href{https://notes.rua.rs/#1/#2.pdf\#nameddest=notes.#3.#4}{#5}}
\newcommand{\missingprereq}[1]{\textcolor{purple}{#1}}
\usepackage{qm-notations}
\usepackage{tikz} % Draw picture
\usetikzlibrary{patterns,decorations.pathmorphing}
\usepackage{tikz-3dplot} 
\lstset{
    numbers=left, 
    numberstyle= \tiny, 
    keywordstyle= \color{ blue!70},
    commentstyle= \color{red!50!green!50!blue!50}, 
    frame=shadowbox, 
    rulesepcolor= \color{ red!20!green!20!blue!20} ,
    escapeinside=``, 
    xleftmargin=2em,xrightmargin=2em, aboveskip=1em,
    framexleftmargin=2em
}  

\title{Mathematical Methods for Waves }
\author{Ruan Haozhe }
\date{September 2025}

\begin{document} 
\setstretch{1}

\clearpage 

\section{Intuition of Quantum Mechanics } 

\subsection{Wave \& Wave Function } 

\textbf{Wave } 
\begin{align}
    \lambda k = 2\pi \quad \omega T = 2\pi 
\end{align} 

\textbf{De Broglie Relation: }
\begin{align}
    p = \frac{h}{\lambda} \quad \hbar = \frac{h}{2\pi} \implies p = \hbar k \quad \hat{p} = -i\hbar \frac{\partial}{\partial x}
\end{align}

\textbf{Schrodinger Equation: } ($\Psi$ is the wave function)
\begin{align}
    i\hbar \frac{\partial}{\partial t} \Psi(x,t) = -\frac{\hbar^2}{2m} \frac{\partial^2}{\partial x^2} \Psi(x,t) + V(x) \Psi(x,t)
\end{align} 

\subsection{Complex Number } 

\textbf{Euler's Formula: } 
\begin{align}
    e^{a + ib} = e^a(\cos b + i \sin b), \quad a \in \mathbb{R}, b \in \mathbb{R} 
\end{align}

\textbf{Complex Conjugate: }
\begin{align}
    z = a + ib \quad z^* = a - ib \quad |z|^2 = z z^* = a^2 + b^2, \quad a,b \in \mathbb{R} \\
    c = e^{\ln(c)} \implies (c^{a+ib})^* = (c^*)^{a-ib}, \quad C \in \mathbb{C}, a,b \in \mathbb{R}
\end{align}

\textbf{Trig $\leftrightarrow$ Hyperbolic: } ($a \in \mathbb{C}$)
\begin{align}
    \sin(ia) = i\sinh(a) \implies \sinh(a) = -i\sin(ia) \\
    \cos(ia) = \cosh(a) \implies \cosh(a) = \cos(ia)
\end{align}

\subsection{Probability } 

\textbf{Gaussian Integral: } ($a > 0$)
\begin{align}
    \int_{-\infty}^{\infty} e^{-ax^2}\,dx = \sqrt{\frac{\pi}{a}}, \quad
    \int_{-\infty}^{\infty} x^2 e^{-ax^2}\,dx = \frac{1}{2a}\sqrt{\frac{\pi}{a}}, \quad
    \int_{-\infty}^{\infty} x^{2n+1} e^{-ax^2}\,dx = 0
\end{align}

\textbf{Axiom of Probability: }
\begin{align}
    \rho(x,t) = |\Psi(x,t)|^2 \quad \int_{-\infty}^{\infty} \rho(x,t) dx = \langle \Psi | \Psi \rangle = 1 \\
    \langle \Psi | \Psi \rangle = C \neq \infty \text{ or } 0 \text{ (normalizable)} \implies \Psi(x,t) \to \frac{1}{\sqrt{C}} \Psi(x,t)
\end{align}

\textbf{Expectation Value: } 
\begin{align}
    \langle O \rangle = \langle \Psi | O \Psi \rangle, \quad \sigma_O = \sqrt{\langle O^2 \rangle - \langle O \rangle^2}, \quad O \in \text{observable}
\end{align} 

\textbf{Commutator: } 
\begin{align}
    \textbf{comm}(A,B) = [A,B] = AB - BA \text{ Use test function}
\end{align}  

\textbf{Uncertainty Principle: }
\begin{align}
    \sigma_A \sigma_B \geq \frac{1}{2} |\langle \textbf{comm}(A,B) \rangle| \implies \sigma_x \sigma_p \geq \frac{\hbar}{2} 
\end{align} 

\subsection{Probability Current } 

\textbf{Probability Density: }
\begin{align}
    \frac{d}{dt} \int^b_a \rho(x,t) dx = J(a,t) - J(b,t) \implies J(x,t) = \frac{\hbar}{2mi} \left( \Psi^* \frac{\partial \Psi}{\partial x} - \Psi \frac{\partial \Psi^*}{\partial x} \right) 
\end{align} 

\textbf{Continuity Equation: }
\begin{align}
    \frac{\partial \rho}{\partial t} + \frac{\partial J}{\partial x} = 0 \text{ (Conservation of Probability)}
\end{align} 

\textbf{Plane Wave } 
\begin{align} 
    \psi(x) = \mathcal{A} e^{ikx} \implies J = \frac{\hbar k}{m} |\mathcal{A}|^2 > 0 \quad \text{(Incident from left)} \\ 
    \psi(x) = \mathcal{B} e^{-ikx} \implies J = -\frac{\hbar k}{m} |\mathcal{B}|^2 < 0 \quad \text{(Incident from right)} 
\end{align} 

\textbf{Transmission \& Reflection Coefficients: }
\begin{align}
    T = \frac{|J_{transmitted}|}{|J_{incident}|}, \quad R = \frac{|J_{reflected}|}{|J_{incident}|} 
\end{align} 

\section{Solving Schrodinger Equation } 

\subsection{TISE } 

\textbf{Separation of Variables: } 
\begin{align}
    \Psi(x,t) = \psi(x) \phi(t) \implies \frac{i\hbar}{\phi(t)} \frac{d\phi(t)}{dt} = \frac{-\frac{\hbar^2}{2m} \frac{d^2 \psi(x)}{dx^2} + V(x) \psi(x)}{\psi(x)} = E 
\end{align} 

\textbf{Time-Dependent Part: } 
\begin{align}
    \phi(t) = e^{-\frac{i}{\hbar} Et} \implies \Psi(x,t) = \psi(x) e^{-\frac{iEt}{\hbar} } 
\end{align} 

\textbf{Time-Independent Schrodinger Equation: }
\begin{align}
    -\frac{\hbar^2}{2m} \frac{d^2}{dx^2} \psi(x) + V(x) \psi(x) = E \psi(x), \quad \hat{H} \psi(x) = E \psi(x) 
\end{align} 

\textbf{Existance of Minium Energy: } 
\begin{align}
    \frac{\hbar^2}{2m} \frac{d^2}{dx^2} \psi = (V(x) - E) \psi, \quad E < V_{min} \implies \text{sign}(\frac{d^2 \psi}{dx^2}) = \text{sign}(\psi) \implies \psi \text{ blows up!} 
\end{align} 

If $E < V_{min}$ the solution is exponentially divergent and not physically acceptable. Therefore allowed energies must be bounded below. 
~ 

\textbf{Scattering State: }
\begin{align}
    \frac{\hbar^2}{2m} \frac{d^2}{dx^2} \psi = (V(x) - E) \psi, \quad E > V_{max} \implies \text{sign}(\frac{d^2 \psi}{dx^2}) \neq \text{sign}(\psi) \implies \psi \text{ oscillates! } 
\end{align} 

$E > V_{max}$ will cause $\psi$ to oscillate, so it is called scattering state. Without linear combination with 2 linear independent states, it is not normalizable. 

\subsection{Infinite Potential Well }

\textbf{Condition: } 
\begin{align}V(x) = \begin{cases}
    0, & 0 < x < a \\
    \infty, & \text{otherwise}
\end{cases} \implies \psi(x) \text{is a $C^0$ function.} 
\end{align} 

\textbf{Equation: } 
\begin{align}
    \lim_{x \to 0^+} \psi(x) = \lim_{x \to 0^-} \psi(x) = \lim_{x \to a^+} \psi(x) = \lim_{x \to a^-} \psi(x) = 0 
\end{align} 

\textbf{Solution: }
\begin{align}
    \psi(x) = A e^{ikx} + B e^{-ikx} \quad \text{where } k = \sqrt{\frac{2mE}{\hbar^2}} \\
    \psi_n(x) = \sqrt{\frac{2}{a}} \sin\left( \frac{n\pi x}{a} \right), \quad E_n = \frac{n^2 \pi^2 \hbar^2}{2ma^2}, \quad n \in \mathbb{N}
\end{align}

\textbf{Expectation Values: }
\begin{align}
    \langle x \rangle = \frac{a}{2}, \quad \langle x^2 \rangle = a^2\!\left(\frac{1}{3} - \frac{1}{2n^2\pi^2}\right), \quad \langle p \rangle = 0, \quad \langle p^2 \rangle = \frac{n^2\pi^2\hbar^2}{a^2}
\end{align}

\subsection{SHM } 

\textbf{Condition: } 
\begin{align}
    V(x) = \frac{1}{2} m \omega^2 x^2 , \quad \omega \in \mathbb{R}^+ \text{ (parameter)} \implies \psi(x) \text{is a $C^\infty$ function.}
\end{align} 

\textbf{Ladder Operators: }
\begin{align}
    a_\pm = \frac{1}{\sqrt{2\hbar m \omega}} \left( \mp i \hat{p} + m \omega x \right) \implies \hat{H} = \hbar \omega \left( a_+ a_- + \frac{1}{2} \right), \textbf{comm}(a_-, a_+) = 1 
\end{align} 

\textbf{Induction: } (We assume $\psi_n(x)$ is an solution - base case. Similar for $a_-$)
\begin{align}
    \hat{H} a_+ \psi_n(x) = a_+ \left( \hat{H} \psi_n(x) + \hbar \omega \psi_n(x) \right) = (E_n + \hbar \omega) a_+ \psi_n(x) 
\end{align} 

\textbf{Ground State: }
\begin{align}
    a_- \psi_0(x) = 0 \implies \psi_0(x) = \left( \frac{m \omega}{\pi \hbar} \right)^{\frac{1}{4}} e^{-\frac{m \omega x^2}{2\hbar}} \\
    a_+ \psi_n(x) = \sqrt{n+1} \psi_{n+1}(x) \quad a_- \psi_n(x) = \sqrt{n} \psi_{n-1}(x) 
\end{align} 

\textbf{General Solution: } 
\begin{align}
    \psi_n(x) = \frac{1}{\sqrt{n!}} (a_+)^n \psi_0(x) \quad E_n = \hbar \omega \left( n + \frac{1}{2} \right)
\end{align}

\textbf{Position \& Momentum: }
\begin{align}
    x = \sqrt{\frac{\hbar}{2m\omega}}\,(a_+ + a_-), \quad \hat{p} = i\sqrt{\frac{\hbar m\omega}{2}}\,(a_+ - a_-)
\end{align}

\textbf{Expectation Values: }
\begin{align}
    \langle x \rangle = 0, \quad \langle x^2 \rangle = \frac{\hbar}{2m\omega}(2n+1), \quad \langle p \rangle = 0, \quad \langle p^2 \rangle = \frac{\hbar m\omega}{2}(2n+1)
\end{align}

\subsection{Free Particle } 

\textbf{Condition: } 
\begin{align}
    V(x) = 0 \implies \frac{d^2}{dx^2} \psi(x) + k^2 \psi(x) = 0, \quad k = \sqrt{\frac{2mE}{\hbar^2} } \in \mathbb{R} 
\end{align} 

\textbf{General Solution: }
\begin{align}
    \psi(x) = A e^{ikx} + B e^{-ikx} \implies \psi(x) = C \sin(kx) + D \cos(kx) 
\end{align} 

\section{Continuous and Fourier Transform }

\subsection{Fourier Transform } 

\textbf{Definition: }
\begin{align}
    \hat{f}(k) = \frac{1}{\sqrt{2\pi}} \int_{-\infty}^{\infty} f(x) e^{-ikx} dx \Longleftrightarrow f(x) = \frac{1}{\sqrt{2\pi}} \int_{-\infty}^{\infty} \hat{f}(k) e^{ikx} dk 
\end{align} 

\subsection{Finite Potential Well } 

\textbf{Condition: }
\begin{align}
    V(x) = \begin{cases}
    0, & 0 < x < a \\
    V_0, & \text{otherwise}
\end{cases} \implies \psi(x) \text{is a $C^1$ function.} 
\end{align} 

\textbf{Equation: } (4 boundary conditions)
\begin{align}
    \forall c, c \in \mathbb{R}, \quad \lim_{x \to c^+} \psi(x) = \lim_{x \to c^-} \psi(x), \quad \lim_{x \to c^+} \frac{d\psi}{dx}(x) = \lim_{x \to c^-} \frac{d\psi}{dx}(x) 
\end{align} 

\textbf{General Solution: } (Scattering, $E > V_0$, wave from left, $B_3 = 0$)
\begin{align}
    \psi(x) = \begin{cases}
    A_1 e^{ik_1 x} + B_1 e^{-ik_1 x}, & x < 0 \\
    A_2 e^{ik_2 x} + B_2 e^{-ik_2 x}, & 0 < x < a \\
    A_3 e^{ik_1 x}, & x > a
\end{cases} \quad k_1 = \sqrt{\frac{2m(E-V_0)}{\hbar^2}}, \quad k_2 = \sqrt{\frac{2mE}{\hbar^2}}
\end{align}

\textbf{Solutions: } (Let $\Delta = 2\cos(k_2 a) - i\left(\dfrac{k_1}{k_2}+\dfrac{k_2}{k_1}\right)\sin(k_2 a)$)
\begin{align}
    A_3 &= \frac{2e^{-ik_1 a}}{\Delta}\,A_1, \quad
    B_1 = \frac{-i\!\left(\tfrac{k_1}{k_2}-\tfrac{k_2}{k_1}\right)\sin(k_2 a)}{\Delta}\,A_1 \\
    A_2 &= \frac{(1+\tfrac{k_1}{k_2})\,e^{-ik_2 a}}{\Delta}\,A_1, \quad
    B_2 = \frac{(1-\tfrac{k_1}{k_2})\,e^{ik_2 a}}{\Delta}\,A_1 \\
    T &= \frac{4}{|\Delta|^2} = \left[1 + \frac{V_0^2 \sin^2(k_2 a)}{4E(E-V_0)}\right]^{-1}, \quad R = 1 - T
\end{align}

\subsection{Finite Potential Barrier }

\textbf{Condition: }
\begin{align}
    V(x) = \begin{cases}
    V_0, & 0 < x < a \\
    0, & \text{otherwise}
\end{cases} \implies \psi(x) \text{is a $C^1$ function.}
\end{align}

\textbf{General Solution: } (Scattering, $E < V_0$ tunneling, wave from left, $B_3 = 0$)
\begin{align}
    \psi(x) = \begin{cases}
    A_1 e^{ikx} + B_1 e^{-ikx}, & x < 0 \\
    A_2 e^{\kappa x} + B_2 e^{-\kappa x}, & 0 < x < a \\
    A_3 e^{ikx}, & x > a
\end{cases} \quad k = \sqrt{\frac{2mE}{\hbar^2}}, \quad \kappa = \sqrt{\frac{2m(V_0-E)}{\hbar^2}}
\end{align}

\textbf{Solutions} ($E < V_0$, let $\Delta = 2\cosh(\kappa a) - i\!\left(\dfrac{k}{\kappa}-\dfrac{\kappa}{k}\right)\sinh(\kappa a)$):
\begin{align}
    A_3 &= \frac{2e^{-ika}}{\Delta}\,A_1, \quad
    B_1 = \frac{-i\!\left(\tfrac{k}{\kappa}+\tfrac{\kappa}{k}\right)\sinh(\kappa a)}{\Delta}\,A_1 \\
    A_2 &= \frac{(1+\tfrac{ik}{\kappa})\,e^{-\kappa a}}{\Delta}\,A_1, \quad
    B_2 = \frac{(1-\tfrac{ik}{\kappa})\,e^{\kappa a}}{\Delta}\,A_1 \\
    T &= \frac{4}{|\Delta|^2} = \left[1 + \frac{V_0^2 \sinh^2(\kappa a)}{4E(V_0-E)}\right]^{-1}, \quad R = 1-T
\end{align}
\textbf{Solutions} ($E > V_0$, replace $\kappa \to -ik'$, i.e.\ $\sinh \to i\sin$, $\cosh \to \cos$):
\begin{align}
    T = \left[1 + \frac{V_0^2 \sin^2(k' a)}{4E(E-V_0)}\right]^{-1}, \quad R = 1-T, \quad k' = \sqrt{\frac{2m(E-V_0)}{\hbar^2}}
\end{align}

\subsection{Delta Potential Well}

\textbf{Condition: }
\begin{align}
    V(x) = -\alpha \delta(x) \implies \psi(x) \text{is a $C^0$ function.} 
\end{align} 

\textbf{Equation: } 
\begin{align}
    \lim_{x \to 0^+} \psi(x) &= \lim_{x \to 0^-} \psi(x) \\ 
    \text{Take }\int_{0^-}^{0^+} dx \text{ for TISE} &\implies \lim_{x \to 0^+} \frac{d\psi}{dx}(x) - \lim_{x \to 0^-} \frac{d\psi}{dx}(x) = -\frac{2m\alpha}{\hbar^2} \psi(0) 
\end{align} 

\textbf{Bound State: } ($E < 0$, one bound state)
\begin{align}
    \psi(x) = \begin{cases}
    A_1 e^{\kappa x}, & x < 0 \\
    B_1 e^{-\kappa x}, & x > 0
\end{cases} \quad \kappa = \frac{m\alpha}{\hbar^2}, \quad E = -\frac{m\alpha^2}{2\hbar^2}, \quad A_1 = B_1 = \sqrt{\kappa}
\end{align}

\textbf{Expectation Values: }
\begin{align}
    \langle x \rangle = 0, \quad \langle x^2 \rangle = \frac{1}{2\kappa^2}, \quad \langle p \rangle = 0, \quad \langle p^2 \rangle = \hbar^2\kappa^2 = -2mE
\end{align}

\textbf{Scattering: } ($E > 0$, wave from left)
\begin{align}
    \psi(x) &= \begin{cases}
    A_1 e^{ikx} + B_1 e^{-ikx}, & x < 0 \\
    A_2 e^{ikx}, & x > 0
\end{cases} \quad k = \sqrt{\frac{2mE}{\hbar^2}}, \quad \beta = \frac{m\alpha}{\hbar^2 k} \\
B_1 &= \frac{i\beta}{1-i\beta} A_1, \quad A_2 = \frac{1}{1-i\beta} A_1, \quad T = \frac{1}{1+\beta^2}, \quad R = \frac{\beta^2}{1+\beta^2}
\end{align}

\subsection{Delta Potential Barrier}

\textbf{Condition: }
\begin{align}
    V(x) = +\alpha \delta(x) \implies \psi(x) \text{ is a $C^0$ function.}
\end{align}

\textbf{Equation: }
\begin{align}
    \lim_{x \to 0^+} \psi(x) = \lim_{x \to 0^-} \psi(x), \quad \lim_{x \to 0^+} \frac{d\psi}{dx}(x) - \lim_{x \to 0^-} \frac{d\psi}{dx}(x) = +\frac{2m\alpha}{\hbar^2} \psi(0)
\end{align}

\textbf{Scattering: } ($E > 0$, wave from left, no bound states)
\begin{align}
    \psi(x) &= \begin{cases}
    A_1 e^{ikx} + B_1 e^{-ikx}, & x < 0 \\
    A_2 e^{ikx}, & x > 0
\end{cases} \quad k = \sqrt{\frac{2mE}{\hbar^2}}, \quad \beta = \frac{m\alpha}{\hbar^2 k} \\
B_1 &= \frac{-i\beta}{1+i\beta} A_1, \quad A_2 = \frac{1}{1+i\beta} A_1, \quad T = \frac{1}{1+\beta^2}, \quad R = \frac{\beta^2}{1+\beta^2}
\end{align}

\clearpage
\section{WKB Approximation (Chapter 8)}

\subsection{WKB Results}

\textbf{Regions and momentum functions:}
\[
  \begin{aligned}
    E>V(x)&:\quad p(x)=\sqrt{2m\cdot(E-V(x))}\quad\text{(allowed)},\\
    E<V(x)&:\quad q(x)=\sqrt{2m\cdot(V(x)-E)}\quad\text{(forbidden)}
  \end{aligned}
\]
The WKB forms below assume a slowly varying potential away from turning points

\textbf{\xref[Allowed-region WKB form]{quantum-mechanics::ch08::thm:wkb-allowed-region}:}
For an allowed interval $I$, fixed $x_0\in I$, and $A,B\in\CC$,
\[
  \psi\approx x\map\frac1{\sqrt{p(x)}}\cdot\left[
    A\cdot\exp\left(\frac{i}{\hbar}\cdot
      \left[\DiffOp^{-1}(p)\right]_{x_0}^{x}\right)
    +B\cdot\exp\left(-\frac{i}{\hbar}\cdot
      \left[\DiffOp^{-1}(p)\right]_{x_0}^{x}\right)\right]
\]

\textbf{\xref[Forbidden-region WKB form]{quantum-mechanics::ch08::thm:wkb-forbidden-region}:}
For a forbidden interval $I$, fixed $x_0\in I$, and $C,D\in\CC$,
\[
  \psi\approx x\map\frac1{\sqrt{q(x)}}\cdot\left[
    C\cdot\exp\left(\frac1\hbar\cdot
      \left[\DiffOp^{-1}(q)\right]_{x_0}^{x}\right)
    +D\cdot\exp\left(-\frac1\hbar\cdot
      \left[\DiffOp^{-1}(q)\right]_{x_0}^{x}\right)\right]
\]
Neglect the growing term only when the physical boundary or matching conditions justify it

\textbf{\xref[Quantization between two infinite walls]{quantum-mechanics::ch08::thm:hard-wall-quantization}:}
For walls at $0,a$ with an allowed interior,
\[
  \psi(0)=\psi(a)=0
  \quad\Longrightarrow\quad
  \left[\DiffOp^{-1}(p)\right]_0^a=n\pi\hbar,
  \qquad n\in\NNOne
\]

\textbf{\xref[Leading WKB tunneling probability]{quantum-mechanics::ch08::thm:wkb-transmission}:}
For a forbidden barrier between $x_L$ and $x_R$,
\[
  \gamma\defeq\frac1\hbar\cdot\left[\DiffOp^{-1}(q)\right]_{x_L}^{x_R},
  \qquad T\approx e^{-2\gamma}\qquad(\gamma\gg1)
\]
For a rectangular barrier of width $L$, $\kappa=\sqrt{2m\cdot(V_0-E)}/\hbar$ gives
\[
  T_{\mathrm{WKB}}\approx e^{-2\kappa L},\qquad
  T_{\mathrm{exact}}\approx\frac{16E\cdot(V_0-E)}{V_0^2}\cdot e^{-2\kappa L}
  \quad(\kappa L\gg1)
\]
The leading exponential agrees; the abrupt barrier edges give a prefactor that WKB's leading exponential omits

\textbf{Gamow's alpha-decay estimate:}
For escape attempts at rate $v/(2R)$ from a nuclear well of radius $R$,
\[
  \tau\approx\frac{2R}{v\cdot T}
  \approx\frac{2R}{v}\cdot e^{2\gamma}
\]

\newpage
\subsection{Airy Functions}

\textbf{\xref[Airy equation and general solution]{quantum-mechanics::ch08::thm:airy-equation-solution}:}
\[
  \DiffOp^2(y)(z)=z\cdot y(z),\qquad
  y:\RR\to\CC\defeq z\map
    a\cdot\AiryAi(z)+b\cdot\AiryBi(z),\quad a,b\in\CC
\]

\textbf{Integral representations:}
The endpoint $\infty$ denotes an improper integral
\[
  \AiryAi:\RR\to\RR\defeq z\map
  \frac1\pi\cdot\left[\DiffOp^{-1}\left(s\map
    \cos\left(\frac{s^3}{3}+s\cdot z\right)\right)\right]_0^\infty
\]
\[
  \AiryBi:\RR\to\RR\defeq z\map
  \frac1\pi\cdot\left[\DiffOp^{-1}\left(s\map
    e^{-s^3/3+s\cdot z}+\sin\left(\frac{s^3}{3}+s\cdot z\right)
    \right)\right]_0^\infty
\]

\subsection{Approximate Forms and Turning-Point Connections}

\textbf{\xref[Airy asymptotic forms]{quantum-mechanics::ch08::thm:airy-asymptotics}:}
For $z\gg0$,
\[
  \begin{aligned}
    \AiryAi(z)&\approx\frac{e^{-2z^{3/2}/3}}{2\sqrt\pi\cdot z^{1/4}},\\
    \AiryBi(z)&\approx\frac{e^{2z^{3/2}/3}}{\sqrt\pi\cdot z^{1/4}}
  \end{aligned}
\]
For $z\ll0$,
\[
  \begin{aligned}
    \AiryAi(z)&\approx\frac1{\sqrt\pi\cdot(-z)^{1/4}}\cdot
      \sin\left(\frac23\cdot(-z)^{3/2}+\frac\pi4\right),\\
    \AiryBi(z)&\approx\frac1{\sqrt\pi\cdot(-z)^{1/4}}\cdot
      \cos\left(\frac23\cdot(-z)^{3/2}+\frac\pi4\right)
  \end{aligned}
\]

\textbf{Local Airy patch:}
At a simple turning point $x_t$, let $s=\DiffOp(V)(x_t)\ne0$ and $V(x_t)=E$
\[
  V(x)\approx E+s\cdot(x-x_t),\qquad
  \alpha\defeq\left(\frac{2m\cdot|s|}{\hbar^2}\right)^{1/3},\qquad
  z=\operatorname{sgn}(s)\cdot\alpha\cdot(x-x_t)
\]
\[
  \psi_P\approx x\map
    a\cdot\AiryAi\bigl(\operatorname{sgn}(s)\cdot\alpha\cdot(x-x_t)\bigr)
    +b\cdot\AiryBi\bigl(\operatorname{sgn}(s)\cdot\alpha\cdot(x-x_t)\bigr)
\]
For a decaying forbidden-side solution of amplitude $D$, matching gives
\[
  a=\frac{2\sqrt\pi\cdot D}{\sqrt{\hbar\cdot\alpha}},\qquad b=0
\]

\newpage
\textbf{\xref[Upward-sloping connection]{quantum-mechanics::ch08::thm:upward-connection-formula}:}
$s>0$: left allowed, right forbidden
\[
  \psi\approx x\map\begin{cases}
    \dfrac{2D}{\sqrt{p(x)}}\cdot\sin\left(
      \dfrac1\hbar\cdot\left[\DiffOp^{-1}(p)\right]_{x}^{x_t}
      +\dfrac\pi4\right),&x<x_t,\\[1em]
    \dfrac{D}{\sqrt{q(x)}}\cdot\exp\left(
      -\dfrac1\hbar\cdot\left[\DiffOp^{-1}(q)\right]_{x_t}^{x}\right),&x>x_t
  \end{cases}
\]

\textbf{\xref[Downward-sloping connection]{quantum-mechanics::ch08::thm:downward-connection-formula}:}
$s<0$: left forbidden, right allowed
\[
  \psi\approx x\map\begin{cases}
    \dfrac{D}{\sqrt{q(x)}}\cdot\exp\left(
      -\dfrac1\hbar\cdot\left[\DiffOp^{-1}(q)\right]_{x}^{x_t}\right),&x<x_t,\\[1em]
    \dfrac{2D}{\sqrt{p(x)}}\cdot\sin\left(
      \dfrac1\hbar\cdot\left[\DiffOp^{-1}(p)\right]_{x_t}^{x}
      +\dfrac\pi4\right),&x>x_t
  \end{cases}
\]
These WKB expressions apply away from $x_t$; use the Airy patch near $x_t$

\textbf{Equivalent phase convention:}
\[
  \sin\left(\theta+\frac\pi4\right)=\cos\left(\theta-\frac\pi4\right)
\]

\end{document}
